\documentclass[12pt]{article}
\usepackage[utf8]{inputenc}
\usepackage[T1]{fontenc}
\usepackage{hyperref}
\title{Torres Strait Creole Language Resources for K-12 Curriculum Development}
\author{Compiled by Claude Sonnet 5 (claude-sonnet-5) for Andr\'as Kornai}
\date{}
\begin{document}
\maketitle
\section{Introduction}
Torres Strait Creole (Yumplatok, Broken) is an English-based creole of the Torres Strait Islands and northern Queensland, Australia. It is the main everyday language of the Islanders; estimates range from about 6,000 to 40,000 users depending on the source and definition. Glottocode torr1261.
\section{Primary Educational Resources}
No dedicated Yumplatok textbook or official curriculum was found. A Queensland Studies Authority profile describes teacher Kali Ahmat at Tagai State College, who applies a language lens to the Australian curriculum so that students write accurately in both English and Yumplatok: \url{https://www.stories.qct.edu.au/post/kali-ahmat}
\section{Linguistic and Reference Materials}
\begin{itemize}
\item Shnukal's report for Thursday Island State High School on differences between Standard Australian English and Torres Strait Creole (found via the Australian Journal of Indigenous Education): \url{https://ajie.atsis.uq.edu.au/ajie/article/download/359/299/918}
\item Another AJIE article found in the same search (content not confirmed): \url{https://ajie.atsis.uq.edu.au/ajie/article/download/376/316/952}
\item Overview of Torres Strait Island languages, including the Ray vocabularies republished by Edwards (these cover the traditional languages, not the creole): \url{https://en.wikipedia.org/wiki/Torres_Strait_Island_languages}
\item Humanities Australia article on Kriol, Yumplatok and Aboriginal English: \url{https://humanities.org.au/power-of-the-humanities/kriol-yumplatok-and-aboriginal-english-australias-contact-languages/}
\item James Cook University repository record: \url{https://researchonline.jcu.edu.au/32593/}
\end{itemize}
\section{Religious and Community Resources}
None identified in the searches.
\section{Digital and Audio Resources}
None identified in the searches.
\section{Curriculum Development Context}
Yumplatok is widely recognised as a language in its own right and is used in some island classrooms, but no standard published curriculum or dictionary was found. Research on transfer from the creole into written English is relevant to bilingual design.
\end{document}